%% ----- develop commands -----------------------------------------------------
\newcommand{\dummyText}[2]{%
    \ifthenelse{\boolean{isDummyText}}{
            {\color{#1}\lipsum[#2]}
    }{}
}

%% ----- Abkürzungsverzeichnis ------------------------------------------------
\makeatletter %% erlaubt (at) symbol as text
\newcommand{\@acronym@nolist}{0}
\AtBeginDocument{
    \@ifpackagewith{acronym}{nolist}
    {\renewcommand{\@acronym@nolist}{1}}
    {\renewcommand{\@acronym@nolist}{0}}
}
%% \listofabbreviations ; add to toc when [intoc], no list (nor title) when package 'acronym' has [nolist]
\newcommand{\listofabbreviations}[1][]{
    \ifthenelse{ \@acronym@nolist=0 }{
        \chapter*{\abbrvTitle} %% add sectoin without numbering (and not in TOC)
        \label{toc:\abbrvTitle}
        \ifthenelse{ \equal{#1}{intoc} }{
            \addcontentsline{toc}{chapter}{\abbrvTitle}
        }{}
    }{}
}
\makeatother %% reset (at) symbol to default behavior


%% ----- Referenzen (außerhalb des Dokument) ----------------------------------
\NewDocumentCommand{\Zitat}{ O{} O{} m }{\parencite[#1][#2]{#3}\xspace}
\NewDocumentCommand{\Vgl}{ O{} m }{\parencite[vgl.][#1]{#2}\xspace}
\newcommand{\Note}[1]{\footnote{#1}}


%% ----- Styling --------------------------------------------------------------
\newcommand{\itemd}[2]{\item{\textbf{#1}}\\{#2}} %% erzeugt ein Listenelement mit fetter Überschrift
\newcommand{\changefont}[3] %% einfaches Wechseln der Schrift, z.B.: \changefont{cmss}{sbc}{n}
{\fontfamily{#1} \fontseries{#2} \fontshape{#3} \selectfont}
\newenvironment{absolutelynopagebreak}
{\par\nobreak\vfil\penalty0\vfilneg
\vtop\bgroup}
{\par\xdef\tpd{\the\prevdepth}\egroup
\prevdepth=\tpd}


%% ----- Tables ---------------------------------------------------------------
\newcommand{\hTable}[4][!htb]{%
%% Optional: placement by h=here t='top' b='bottom' p='extra page', enforce with '!'
%% Required: caption (in text), label (for reference), table env
    \begin{table}[#1]
        \begin{threeparttable}[b]
            \small
            \onehalfspacing
            \setlength{\abovecaptionskip}{0em plus 2em minus 2em}
            \caption{#2}\label{#3}
            \begin{center}
                #4
            \end{center}
        \end{threeparttable}
    \end{table}
}
\newcommand{\vTable}[3]{%
    \begin{landscape}
        \hTable{#1}{#2}{#3}
    \end{landscape}
}

\newcommand{\vCell}[2][]{% %% Optional: text lines to span/center text
    \ifthenelse{\isempty{#1}}{%
        \vspace*{-.65em}\rotatebox[origin=c]{90}{#2}\vspace*{.3em}%
    }{%
        \multirow{#1}{*}{\rotatebox{90}{#2}}%
    }
%\rotatebox{90}{#1}%
%\rotatebox[origin=c]{90}{\parbox[t]{1em}{#1}}%
%\rotatebox[origin=c]{90}{\centering\hfill#1\hfill}%
%\rotatebox{90}{\centering\hfill#1\hfill}%
%\rotatebox{90}{\hfill#1\hfill}%
%\parbox[t]{2mm}{\multirow{3}{*}{\rotatebox[origin=c]{#1}}}%
%\begin{tabular}{@{}c@{}}\rotatebox[origin=c]{#1}\end{tabular}&
}


%% ----- Graphics -------------------------------------------------------------
\newcommand{\Image}[2][]{%
%% usage like \includegraphics : \Image[<includegraphics-options>]{<path>}
    \includegraphics[#1]{#2}%
}
\newcommand{\ImageInRatio}[3][]{%
%% usage: \ImageInRatio[<includegraphics-options>]{<scale>}{<path>}
    \Image[width=#2\textwidth,height=#2\textheight,keepaspectratio,#1]{#3}%
}


%% ----- Semantische Textverarbeitung -----------------------------------------
\makeatletter
\let\ACplacelabel\AC@placelabel %% required for hyperref, else warning about undefined ref
\makeatother

\newcommand{\Abk}[1]{\textit{\ac{#1}}} % Abkürzung
\newcommand{\AbkK}[1]{\textit{\acs{#1}}\ACplacelabel{#1}} % Abkürzung kurz (ohne Erstnennung ausgeschrieben)
\newcommand{\AbkL}[1]{\textit{\acf{#1}}\ACplacelabel{#1}} % Abkürzung lang (immer ausgeschrieben)
\newcommand{\AbkS}[2]{#2 (#1)} % Abkürzung: one-time (single) = nicht im Verzeichnis
\newcommand{\Datum}[1]{\DTMdate{#1}} % Datum
\newcommand{\Monat}[1]{% Datum ohne Tag, z.B. \Monat{2020-09}
    \DTMlangsetup{showdayofmonth=false}%
    \DTMdate{#1-01}%
    \DTMlangsetup{showdayofmonth=true}%
}
\newcommand{\Jahr}[1]{#1} % Datum nur als Jahr
\newcommand{\Englisch}[1]{ (englisch: \textit{#1})} % englischer Begriff
\newcommand{\Begriff}[1]{\textit{#1}}
\newcommand{\Tech}[1]{\textsf{#1}}

\newcommand{\Klasse}[1]{\texttt{#1}}
\newcommand{\Methode}[1]{\texttt{#1}}
\newcommand{\MethodenRef}[1]{\texttt{#1}}
\newcommand{\Code}[1]{\textsf{#1}}

%% ----- text sizes -----------------------------------------------------------
%% \Huge
%% \huge
%% \LARGE
%% \Large
%% \large
%% \normalsize (default)
%% \small
%% \footnotesize
%% \scriptsize
%% \tiny


%% ----- Verweise (im Dokument) -----------------------------------------------
%% Default  = Nr, Name
%% K : Kurz = Nr, Seite
%% V : Voll = Nr, Name, Seite
%% https://tex.stackexchange.com/questions/149807/autoref-subsections-in-appendix

%% automatic naming with \autoref, using custom command for better readability in source
\newcommand{\Gleichung}[1]{\autoref{#1}}
\newcommand{\Abbildung}[1]{\autoref{#1}}
\newcommand{\Anhang}[1]{\autoref{#1}}

\newcommand{\TabelleDefaultRef}{KNP}
\ifthenelse{\boolean{isBookPrint}}{}{
    \renewcommand{\TabelleDefaultRef}{KN}
}
\newcommand{\Tabelle}[2][\TabelleDefaultRef]{% K=Tabelle~<Nr> N=<Tabellenname> P='auf S.~<Seite>'
    \IfStrEqCase{#1}{
            {N}{\nameref{#2}}
            {NK}{\hyperref[#2]{\nameref{#2} (\autoref{#2})}}
            {KN}{\hyperref[#2]{\autoref{#2}:~\nameref{#2}}}
            {KP}{\hyperref[#2]{\autoref{#2} auf S.~\pageref{#2}}}
            {KNP}{\hyperref[#2]{\autoref{#2}:~\nameref{#2} auf S.~\pageref{#2}}}
    }[\autoref{#2}]%
}

\newcommand{\KapitelDefaultRef}{KNP}
\ifthenelse{\boolean{isBookPrint}}{}{
    \renewcommand{\KapitelDefaultRef}{KN}
}
\newcommand{\Kapitel}[2][\KapitelDefaultRef]{% K=Kapitel~<Nr> N=<Kapitelname> P='auf S.~<Seite>'
    \IfStrEqCase{#1}{
            {N}{\nameref{#2}}
            {NK}{\hyperref[#2]{\nameref{#2} (\autoref{#2})}}
            {KN}{\hyperref[#2]{\autoref{#2}:~\nameref{#2}}}
            {KP}{\hyperref[#2]{\autoref{#2} auf S.~\pageref{#2}}}
            {KNP}{\hyperref[#2]{\autoref{#2}:~\nameref{#2} auf S.~\pageref{#2}}}
    }[\autoref{#2}]%
}

%% manual naming with \ref
%\newcommand{\Gleichung}[1]{Gleichung~\ref{#1}}
%\newcommand{\Abbildung}[1]{Abbildung~\ref{#1}}
%\newcommand{\AbbildungV}[1]{Abbildung~\ref{#1}:~\nameref{#1} auf S.~\pageref{#1}}
%\newcommand{\Anhang}[1]{Anhang~\ref{#1}:~\nameref{#1} auf S.~\pageref{#1}}
%\newcommand{\AnhangK}[1]{Anhang~\ref{#1} auf S.~\pageref{#1}}
%\newcommand{\AnhangKapitel}[1]{Kapitel~\ref{#1}:~\nameref{#1} auf S.~\pageref{#1}}
%\newcommand{\Tabelle}[1]{Tabelle~\ref{#1}:~\nameref{#1}}
%\newcommand{\TabelleK}[1]{Tabelle~\ref{#1} auf S.~\pageref{#1}}
%\newcommand{\TabelleV}[1]{Tabelle~\ref{#1}:~\nameref{#1} auf S.~\pageref{#1}}
%\newcommand{\Verweis}[1]{Kapitel~\ref{#1}}
%\newcommand{\VerweisKL}[1]{Kapitel~\ref{#1} auf Seite~\pageref{#1}}
%\newcommand{\VerweisVL}[1]{Kapitel~\ref{#1}:~\nameref{#1} auf Seite~\pageref{#1}}
%\newcommand{\VerweisKT}[1]{(Siehe Kapitel~\ref{#1} auf Seite~\pageref{#1})}
%\newcommand{\VerweisVT}[1]{(Siehe Kapitel~\ref{#1}:~\nameref{#1} auf Seite~\pageref{#1})}
%\newcommand{\VerweisKF}[1]{\footnote{Siehe Kapitel~\ref{#1} auf Seite~\pageref{#1}.}}
%\newcommand{\VerweisVF}[1]{\footnote{Siehe Kapitel~\ref{#1}:~\nameref{#1} auf Seite~\pageref{#1}.}}


%% ----- Abkürzungen, ggf. mit korrektem Leerraum ----------------------------
\newcommand{\bs}{$\backslash$\xspace}
\newcommand{\bspw}{bspw.\xspace} % beispielsweiße
\newcommand{\bzw}{bzw.\xspace} % beziehungsweiße
\newcommand{\ca}{ca.\xspace} % circa
\newcommand{\dahe}{\mbox{d.\,h.}\xspace} % das heißt
\newcommand{\etc}{etc.\xspace} % etcetera (und so weiter (usw))
\newcommand{\eur}[1]{\mbox{#1\,\texteuro}\xspace} % Euro Zeichen
\newcommand{\evtl}{evtl.\xspace} % eventuell
\newcommand{\ggf}{ggf.\xspace} % gegebenenfalls
\newcommand{\Ggf}{Ggf.\xspace} % Gegebenenfalls
\newcommand{\gqq}[1]{\glqq{}#1\grqq{}}
\newcommand{\inkl}{inkl.\xspace} % inklusive
\newcommand{\insb}{insb.\xspace} % insbesondere
\newcommand{\uA}{\mbox{u.\,A.}\xspace} % unter Anderem
\newcommand{\usw}{usw.\xspace} % und so weiter
\newcommand{\vgl}{vgl.\xspace} % vergleich
\newcommand{\Vvgl}{Vgl.\xspace} % Vergleich
\newcommand{\zB}{\mbox{z.\,B.}\xspace} % zum Beispiel
\newcommand{\ZB}{\mbox{Z.\,B.}\xspace} % zum Beispiel
\newcommand{\og}{\mbox{o.\,g.}\xspace} % oben genannt
\newcommand{\sog}{sog.\xspace} % sogenannt
